\documentclass[a4paper]{article}
\usepackage{amsmath}
\usepackage{amsfonts}
\usepackage[affil-it]{authblk}
\usepackage[backend=bibtex,style=numeric,sorting=none]{biblatex}
\usepackage{enumitem}
\usepackage{graphicx}
\usepackage{float}

\usepackage{geometry}
\geometry{margin=1.5cm, vmargin={0pt,1cm}}
\setlength{\topmargin}{-1cm}
\setlength{\paperheight}{29.7cm}
\setlength{\textheight}{25.3cm}
\setlength{\parindent}{0pt}

\begin{document}
% =================================================
\title{Numerical Analysis theoretical homework \# 6}

\author{WangHao 3220103613
  \thanks{Electronic address: \texttt{3220103613@zju.edu.cn}}}
\affil{(Information and Computing Sciences, Class 2201), Zhejiang University }


\date{Due time: \today}

\maketitle
% ============================================
\section*{I.}

\subsection*{(a)}

Simpson's rule:
\begin{equation*}
I^S(y) = \frac{(-1)}{6} \left[ f(-1) + 4f(0) + f(1) \right] = \frac{1}{3} f(-1) + \frac{4}{3} f(0) + \frac{1}{3} f(1)
\end{equation*}

\[
\begin{array}{|c|cccc|}
\hline
-1 & y(-1) &  &  & \\
0 & y(0) & y(0)-y(1) & & \\
0 & y(0) & y'(0) & y'(0)-y(0)+y(-1) & \\
1 & y(1) & y(1)-y(0) & y(1)-y(0)-y'(0) & \frac{1}{2}[y(1)-2y'(0)-y(-1)]\\
\hline
\end{array}
\]

\begin{equation*}
\Rightarrow p_3(x) = y(-1) + [y(0) - y(-1)](x + 1) + [y'(0) - y(0) + y(-1)](x + 1)x + \frac{1}{2} [y(1) - 2y'(0) - y(-1)](x + 1)x^2
\end{equation*}

\begin{equation*}
\int_{-1}^{1} p_3(t) \, dt = 2y(-1) + 2[y(0) - y(-1)] + \frac{2}{3} [y'(0) - y(0) + y(-1)] + \frac{2}{3} \cdot \frac{1}{2} [y(1) - 2y'(0) - y(-1)]
\end{equation*}

\begin{equation*}
= \frac{1}{3} y(-1) + \frac{4}{3} y(0) + \frac{1}{3} y(1),
\end{equation*}

which is same as Simpson's rule.

\subsection*{(b)}

From Hermite interpolation formula, we have: 

\begin{equation*}
y(t) = p_3(t) + \frac{y^{(4)}(\xi)}{4!} (t+1)t^2(t-1)
\end{equation*}

\begin{equation*}
\Rightarrow \int_{-1}^{1} y(t) \, dt = \int_{-1}^{1} p_3(t) \, dt + \int_{-1}^{1} \frac{y^{(4)}(\xi)}{4!} (t+1)t^2(t-1) \, dt
\end{equation*}

\begin{equation*}
\Rightarrow E^S(y) = \int_{-1}^{1} \frac{y^{(4)}(\xi)}{4!} (t+1)t^2(t-1) \, dt
\end{equation*}

\begin{equation*}
\text{Since } (t+1)t^2(t-1) \text{ do not change its sign over the interval } [-1, 1], \text{ by the mean value theorem of integrals:}
\end{equation*}

\begin{equation*}
\exists \xi \in (-1,1),\quad \text{s.t.}
\end{equation*}
\begin{equation*}
E^S(y) = \frac{y^{(4)}(\xi)}{4!} \int_{-1}^{1} (t+1)t^2(t-1) \, dt = -\frac{1}{90} y^{(4)}(\xi)
\end{equation*}

\subsection*{(c)}

Assume applying Simpson's rule on $[a,b]$.

Similarly, there exists $\xi \in (a, b)$, such that
\begin{equation*}
E^S(y) = \frac{y^{(4)}(\xi)}{4!} \int_a^b (t-a)\left(t-\frac{a+b}{2}\right)^2(t-b) \, dt = -\frac{y^{(4)}(\xi)}{2880}(b-a)^5
\end{equation*}

Divide $[a,b]$ into $n$ segments, denoted as $0, 1, \ldots, n-1$, with $h = \frac{b-a}{n}$. The points are $a_0, a_1, \ldots, a_{2n-2}, a_{2n-1}, a_{2n}$, where $a_i = a + \frac{1}{2} ih$. The $i$-th segment contains points $a_{2i}, a_{2i+1}, a_{2i+2}$, within the interval $[a+ih, a+(i+1)h]$.

Applying Simpson's rule on each segment, we have:
\begin{equation*}
I(y) = \sum_{i=0}^{n-1} h \frac{1}{6} \left( f(a_{2i}) + 4f(a_{2i+1}) + f(a_{2i+2}) \right)
\end{equation*}

\begin{equation*}
= \frac{1}{6} h \left[ f(a_0) + f(a_{2n+2}) \right] + \frac{5}{6} h \left[ f(a_1) + f(a_{2n+1}) \right] + h \sum_{i=2}^{2n} f(a_i)
\end{equation*}

And,

\begin{equation*}
E(y) = \sum_{i=0}^{n-1} -\frac{y^{(4)}(\xi)}{2880} h^5, \text{ where } \xi \in (a_{2i}, a_{2i+2})
\end{equation*}

\begin{equation*}
\leq -(b-a) \frac{y^{(4)}(\xi)}{2880} h^4, \text{ where } \xi \in (a, b)
\end{equation*}

\section*{II.}

\subsection*{(a)}

Divide the interval into $n$ subintervals.\\
For the trapezoidal rule, with $h = \frac{1}{n}$.

\begin{equation*}
R = -\frac{b-a}{12} h^2 f''\left(\xi\right)
\end{equation*}

\begin{equation*}
= -\frac{1}{12n^2} f''\left(\xi\right), \quad \xi \in (0,1)
\end{equation*}

Given the function:
\begin{equation*}
f(x) = e^{-x^2}, \quad f'(x) = -2x e^{-x^2}, \quad f''(x) = (4x^2 - 2) e^{-x^2}
\end{equation*}

\begin{equation*}
f'''(x) = (-8x^3 + 12x) e^{-x^2} = 0 \Rightarrow x = 0, \pm \sqrt{\frac{3}{2}}
\end{equation*}

So
\[
\max_{x \in [0,1]} |f''(x)| = \max \{ f''(0), f''(1) \} = 2
\]

From
\begin{equation*}
|R| < 0.5 \times 10^{-6}
\end{equation*}

\begin{equation*}
\Rightarrow n > \frac{1}{\sqrt{3}} \times 10^3 = 577.3
\end{equation*}

\begin{equation*}
\text{i.e., } n = 578
\end{equation*}

\subsection*{(b)}

For the Simpson's rule, with $h = \frac{1}{2n}$.

\begin{equation*}
R = -\frac{b-a}{180} h^4 f^{(4)}(\xi), \quad \xi \in (0,1)
\end{equation*}

\begin{equation*}
f^{(4)}(x) = (16x^4 - 48x^2 + 12) e^{-x^2}
\end{equation*}

\begin{equation*}
\Rightarrow \max_{x \in [0,1]} |f^{(4)}(x)| = 12
\end{equation*}

From
\begin{equation*}
|R| < 0.5 \times 10^{-6}
\end{equation*}

\begin{equation*}
\Rightarrow n = 9.55
\end{equation*}

\begin{equation*}
\text{i.e., } n \geq 10
\end{equation*}

\section*{III.}

\subsection*{(a)}

$\text{Let } p(t) = mt + n. \\$
\begin{align*}
&\int_{0}^{+\infty} p(t) \pi_2(t) P(t) = \int_{0}^{+\infty} (mt+n)(t^2+at+b) e^{-t} dt \\
&= \int_{0}^{+\infty} (mt^3 + (ma+n)t^2 + (na+mb)t + nb) e^{-t} dt \\
&= m \cdot 3! + (ma+n) \cdot 2! + (na+mb) \cdot 1! + nb \cdot 0! \\
&= 6m + 2ma + 2n + na + mb + nb \\
&= (6 + 2a + b)m + (2 + a + b)n = 0 \\
&\Rightarrow \begin{cases}
6 + 2a + b = 0 \\
2 + a + b = 0
\end{cases} \\
&\Rightarrow a = -4, \quad b = 2. \\
&\Rightarrow \pi_2(t) = t^2 - 4t + 2
\end{align*}

\subsection*{(b)}

\begin{align*}
t &= \frac{4 \pm \sqrt{16 - 8}}{2} = 2 \pm \sqrt{2}, \\
t_1 &= 2 - \sqrt{2}, \quad t_2 = 2 + \sqrt{2}, \\
\omega_1 &= \int_0^{+\infty} \frac{x - t_2}{t_1 - t_2} e^{-t} \, dt = \frac{1 - t_2}{t_1 - t_2} = \frac{1 + \sqrt{2}}{2\sqrt{2}} = \frac{2 + \sqrt{2}}{4}, \\
\omega_2 &= \int_0^{+\infty} \frac{x - t_1}{t_2 - t_1} e^{-t} \, dt = \frac{1 - t_1}{t_2 - t_1} = \frac{\sqrt{2} - 1}{2\sqrt{2}} = \frac{2 - \sqrt{2}}{4}.
\end{align*}

Consider the Hermite interpolation polynomial \( p_3(t) \) at \( \{t_1, t_1, t_2, t_2\} \):

\begin{align*}
f(t) &= p_3(t) + \frac{f^{(4)}(\xi)}{4!} (t - t_1)^2 (t - t_2)^2, \\
E_2(f) &= \int_0^{+\infty} \frac{f^{(4)}(\xi)}{4!} (t - t_1)^2 (t - t_2)^2 e^{-t} \, dt.
\end{align*}

Since \( (t - t_1)^2 (t - t_2)^2 e^{-t} \) do not change its sign over \( (0, +\infty) \), by the mean value theorem of integrals, \\
$\exists \tau > 0$ such that:

\begin{align*}
E_2(f) &= \frac{f^{(4)}(\tau)}{4!} \int_0^{+\infty} (t - t_1)^2 (t - t_2)^2 e^{-t} \, dt \\
&= \frac{f^{(4)}(\tau)}{4!} \left( 4! + (-2t_1 - 2t_2) 3! + (t_1^2 + t_2^2 + 4t_1 t_2) 2! + (-2t_1 t_2^2 - 2t_2 t_1^2) 1! + (t_1^2 t_2^2) 0! \right) \\
&= \frac{f^{(4)}(\tau)}{6}.
\end{align*}

\subsection*{(c)}
\[
I(f) = \frac{2+\sqrt{2}}{4} \cdot \frac{1}{1+t_1} + \frac{2-\sqrt{2}}{4} \cdot \frac{1}{1+t_2}
\]
\[
= \frac{8+5\sqrt{2}}{28} + \frac{8-5\sqrt{2}}{28}
\]
\[
= \frac{4}{7}
\]
\[
(\approx 0.571428)
\]

\[
E_2(f) = \frac{f^{(4)}(t)}{6}
\]
\[
\text{From } f^{(4)}(t) = 4!(t+1)^{-5}
\]
\[
\Rightarrow E_2(f) = 4(\tau+1)^{-5}, \quad \tau > 0
\]

\[
\text{From} \quad E_2(f) = I - \frac{4}{7} \quad \text{we solve that} \quad \tau = 1.7612556005
\]

\section*{IV.}

\subsection*{(a)}

\begin{align*}
  &\ l_m(x_n) = \delta_{m,n} \\
  &\ \text{And} \quad l_m'(x_m) = \frac{1}{\pi_{n+1}(x_m)} \sum_{j \neq m} \prod_{i \neq m, j} (x_m - x_i) = \sum_{j \neq m} \frac{1}{x_m - x_j}
\end{align*}

\begin{align*}
  &\ h_m(x_n) = (a_m + b_m x_n) \delta_{m,n} \Rightarrow a_m + b_m x_m = 1 \\
  &\ h_m'(x_n) = b_m l_m^2(x_n) + (a_m + b_m x_n) 2 l_m(x_n) l_m'(x_n) \\
  &\ = l_m(x_n) (b_m l_m(x_n) + 2 (a_m + b_m x_n) l_m'(x_n)) \\
  &\ = \delta_{m,n} (b_m l_m(x_n) + 2 (a_m + b_m x_n) l_m'(x_n))
\end{align*}

\begin{align*}
  &\ \text{Only need} \quad h'_m (x_m) = 0 \Rightarrow b_m + 2(a_m + b_m x_m) l'_m (x_m) = 0 \\
  &\ \Rightarrow \left\{ \begin{array}{l} a_m + b_m x_m = 1 \\ 2a_m + \left[2x_m +  \frac{1}{\sum_{j \neq m} \frac{1}{x_m - x_j}}\right] b_m = 0 \end{array} \right. \\
  &\ \Rightarrow \left\{ \begin{array}{l} a_m = 1 + 2x_m \sum_{j \neq m} \frac{1}{x_m - x_j} \\ b_m = -2 \sum_{j \neq m} \frac{1}{x_m - x_j} \end{array} \right.
\end{align*}

\begin{align*}
  &\ q_m(x_n) = (c_m + d_m x_n) \delta_{m,n} \equiv 0 \\
  &\ \implies c_m + d_m x_m = 0 \\
  &\ q_m'(x_n) = \delta_{mn} (d_m l_m(x_m) + 2(c_m + d_m x_m) l_m'(x_n)) = \delta_{mn} \\
  &\ \implies d_m l_m(x_m) + 2(c_m + d_m x_m) l_m'(x_m) = 1 \\
  &\ \left\{\begin{array}{l}c_m + x_m d_m = 0 \\ 2 l_m'(x_m) c_m + (2 x_m l_m'(x_m) + 1) d_m = 1\end{array}\right. \\
  &\ \Rightarrow \left\{\begin{array}{l}c_m = -x_m \\ d_m = 1\end{array}\right.
\end{align*}

\subsection*{(b)}
\begin{align*}
  &\ \text{Perform Hermite interpolation on } p \text{ at } x = x_1, x_1, x_2, x_2, \ldots, x_n, x_n \text{ to obtain } p_{2n+1}(x). \\
  &\ \text{Take } I_n(f) = I(P_{2n-1}(x)) = I\left[\sum_{i=1}^{n} (h_i(x)f_i + g_i(x)f_i')\right] = \sum_{i=1}^{n} [I(h_i(x))f_i + I(g_i(x))f_i'] \\
  &\ \left\{\begin{array}{l}w_k = \int_a^b (a_k + b_k t) l_k^2(t) \rho(t) dt \\ \mu_k = \int_a^b (c_k + d_k t) l_k^2(t) \rho(t) dt\end{array}\right. \\
  &\ \text{For all } p \in \mathbb{P}_{2n-1}, \text{ we have } p_{2n+1}(x) = p \Rightarrow I_n(p) = I(p_{2n-1}) = I(p) \Rightarrow E_n(p) = 0 
\end{align*}

\subsection*{(c)}
\begin{align*}
   \mu_k &= \int_a^b (t - x_k) l_k^2(t) \rho(t) dt \\
  &\ =\frac{1}{\prod_{i \neq k} (x_i - x_k)} \int_a^b n(t) l_k(t) \rho(t) dt = 0 \quad \text{(where } n(t) = \prod_{i=1}^n (t - x_i) \text{)} , \text{For } k = 1, 2, \ldots, n \\
  &\ \text{i.e. }\int_a^b n(t) l_k(t) \rho(t) dt = 0 \quad (k = 1, 2, \ldots, n) \\
  &\ \text{Because } l_1(t), l_2(t), \ldots, l_n(t) \text{ are the basis of } \mathbb{P}_{n-1} \\
  &\ \Rightarrow \text{node polynomial }n(t) \text{ is an } n \text{th order orthogonal polynomial of $\rho(t)$ and $[a,b]$.}
\end{align*}

\end{document}